\Subhead{3.8.1}{Пространственная и временная когерентность. Дифракция случайных волн. Теорема Ван Циттерта--Цернике}

Для комплексного аналитического поля определим $\Gamma_{12}(\tau)=\langle E(\bm r_1,t+\tau)E^*(\bm r_2,t)\rangle$. Нормированная степень когерентности $\gamma_{12}=\Gamma_{12}/\sqrt{I_1I_2}$ удовлетворяет $|\gamma_{12}|\leq1$ по Коши–Буняковскому. При сложении двух пучков средняя интенсивность содержит $2\operatorname{Re}\Gamma_{12}$; видность полос
\begin{equation}
 V=\frac{2\sqrt{I_1I_2}}{I_1+I_2}|\gamma_{12}|.
\end{equation}
Временная когерентность связана со спектральной шириной: $\tau_c\sim1/\Delta\omega$ по порядку, точный коэффициент зависит от формы спектра и определения ширины. Пространственная характеризует корреляцию разных точек в один момент.

Линейное распространение $E(\bm r)=\int h(\bm r,\bm\rho)E_s(\bm\rho)\dd^2\rho$ даёт
\begin{equation}
 \Gamma_{12}=\iint h_1(\bm\rho_1)h_2^*(\bm\rho_2)
 \Gamma_s(\bm\rho_1,\bm\rho_2)\dd^2\rho_1\dd^2\rho_2.
\end{equation}
Для квазимонохроматического пространственно некогерентного источника $\Gamma_s=I_s(\bm\rho_1)\delta(\bm\rho_1-\bm\rho_2)$. В дальней параксиальной области разность фаз до двух точек с базой $\bm b$ линейна по $\bm\rho$, поэтому после удаления известной геометрической фазы
\begin{equation}
 \gamma(\bm b)=\frac{\int I_s(\bm\rho)\exp[-\ii k\bm b\cdot\bm\rho/z]\dd^2\rho}
 {\int I_s(\bm\rho)\dd^2\rho}.
\end{equation}
Это теорема Ван Циттерта–Цернике. Больший угловой размер источника уменьшает размер области когерентности: $b_c\sim\lambda/\theta_s$. Звёздный интерферометр измеряет угловую структуру источника через видность, хотя изображение отдельным телескопом может не разрешаться.
